\section{Notation and Setup}\label{sec:notation}

All divisibility statements in this paper are inside the finite board
\[
  V_n:=\{2,3,\ldots,n\}.
\]
A set $A\subseteq V_n$ is \emph{primitive} if no two distinct elements of $A$
are comparable under divisibility.  Equivalently, $A$ is an antichain in the
divisibility poset.  A primitive set is \emph{maximal} if no element of
$V_n\setminus A$ can be added while preserving primitiveness.

The game value $\L(n)$ is the final size of the primitive set under optimal
play when Prolonger moves first and maximizes the final size, while Shortener
minimizes it.

We use the lower--upper split
\[
  L_n:=\{m\in\mathbb Z:\ 2\le m\le \lfloor n/2\rfloor\},
  \qquad
  \U_n:=\{m\in\mathbb Z:\ n/2<m\le n\}.
\]
The upper half $\U_n$ is an antichain, since two distinct elements of $\U_n$
cannot divide each other.

For $x\in L_n$, the upper shadow and shadow weight are
\[
  M_n(x):=\{u\in\U_n:\ x\mid u\},
  \qquad
  w_n(x):=|M_n(x)|-1.
\]
The subtraction by $1$ is convenient because a lower-half move $x$ adds one
move but may remove $|M_n(x)|$ possible upper-half moves from a terminal
position.

For $P\subseteq\U_n$, define
\[
  \mathcal L_n(P):=\{x\in L_n:\ x\nmid u\ \text{for every }u\in P\},
\]
and
\[
  \beta_n(P):=
  \max\left\{\sum_{x\in B}w_n(x):
    B\subseteq\mathcal L_n(P)\ \text{is a divisibility antichain}\right\}.
\]
The maximum is over a finite set, and is understood to be $0$ when
$\mathcal L_n(P)$ is empty.

We use standard arithmetic notation.  The functions $\pi(x)$ and
$\vartheta(x)$ denote the prime-counting function and Chebyshev's function,
respectively.  For an integer $m$, $P^-(m)$ and $P^+(m)$ denote the least and
greatest prime factors when $m>1$, $\Omega(m)$ denotes the number of prime
factors counted with multiplicity, and $\omega(m)$ denotes the number of
distinct prime factors.  The symbol $\PP$ denotes the set of primes.  As usual,
$O(\cdot)$, $o(\cdot)$, $\ll$, and $\gg$ are interpreted as $n\to\infty$ unless
a different parameter is explicitly specified; subscripts indicate dependence
of the implied constants.
